\documentclass[12pt,a4paper]{article}
\usepackage[utf8]{inputenc}
\usepackage[vietnamese]{babel}
\usepackage{amsmath,amssymb}
\usepackage[margin=2cm]{geometry}
\usepackage{enumitem}
\usepackage{fontspec}
\setmainfont{Times New Roman}

\title{\textbf{BÀI TẬP VỀ NHÀ: BÀI 4 -- NHIỆT DUNG RIÊNG}\\
\large Vật Lý 12 -- Chương trình GEMS}
\author{Giáo viên: Kha Khung Hiệp}
\date{}

\begin{document}
\maketitle
\textit{Thời gian làm bài: 45 phút}

\section*{PHẦN I: CÂU HỎI TRẮC NGHIỆM NHIỀU LỰA CHỌN (18 CÂU)}

\textbf{Câu 1:} Đơn vị đo của nhiệt dung riêng trong hệ SI là:
\begin{enumerate}[label=\Alph*.]
  \item $\text{J/kg}$.
  \item $\text{J/K}$.
  \item $\text{J/kg.K}$.
  \item $\text{W/kg.K}$.
\end{enumerate}

\textbf{Câu 2:} Ý nghĩa của con số nhiệt dung riêng của nhôm $c = 880\text{ J/kg.K}$ là:
\begin{enumerate}[label=\Alph*.]
  \item Để 1 kg nhôm nóng chảy hoàn toàn cần cung cấp nhiệt lượng 880 J.
  \item Để 1 kg nhôm tăng thêm $1^{\circ}\text{C}$ cần cung cấp nhiệt lượng 880 J.
  \item Khi 1 kg nhôm giảm đi 880 K thì nó giải phóng nhiệt lượng 1 J.
  \item Cần nhiệt năng 880 W để đun nóng nhôm lên 1 độ C trong 1 giây.
\end{enumerate}

\textit{(Các câu hỏi trắc nghiệm từ Câu 3 đến Câu 18 về tính toán nhiệt lượng đun nóng chất lỏng, so sánh các chất...)}

\section*{PHẦN II: CÂU HỎI ĐÚNG/SAI (4 CÂU, MỖI CÂU 4 Ý)}

\textbf{Câu 1:} Trong thí nghiệm đo nhiệt dung riêng của nước bằng nhiệt lượng kế:
\begin{enumerate}[label=\alph*)]
  \item Bình nhiệt lượng kế xốp có vỏ nhựa nhằm mục đích cách nhiệt tốt với môi trường bên ngoài. (Đ/S)
  \item Công thức xác định nhiệt dung riêng thực nghiệm là $c = \dfrac{P \cdot m}{t \cdot \Delta t}$. (Đ/S)
  \item Nếu đun nước quá lâu, nhiệt lượng thất thoát ra môi trường nhiều hơn sẽ làm cho giá trị đo được của nhiệt dung riêng nhỏ hơn giá trị thực. (Đ/S)
  \item Dùng máy khuấy đều nước trong bình giúp nhiệt độ phân bố đều, giảm thiểu sai số đo nhiệt độ. (Đ/S)
\end{enumerate}

\section*{PHẦN III: CÂU HỎI TRẢ LỜI NGẮN (6 CÂU)}

\textbf{Câu 1:} Đun nóng một khối kim loại đồng có khối lượng $m = 0{,}5\text{ kg}$ từ nhiệt độ $25^{\circ}\text{C}$ lên $85^{\circ}\text{C}$ cần cung cấp nhiệt lượng $11400\text{ J}$. Hãy tính nhiệt dung riêng của đồng (đơn vị: $\text{J/kg.K}$).

\textit{(Đáp số ghi dưới dạng số nguyên)}

\vfill
\begin{center}
\textit{--- GEMS Vật Lý 12 -- Bài 4: Nhiệt Dung Riêng ---}
\end{center}

\end{document}
